% Bewertungspaket zum Gesamttest «Wärmeausdehnung und Gase» — Quelle des PDFs (python3 scripts/build-lp-pdf.py ausdehnung).
% Ganze Punkte je Teilschritt; Werte und Folgefehler-Fälle mit python3 nachgerechnet (07.10.2026).
% Aufbau, Auftrag an die KI und Punktarten (E, A, W, B) wie im Leitprogramm Hydrostatik; Zeilenabstand im
% Raster 1.5, «ein Fehler, ein Abzug» auch über Teilaufgaben, Rundung 2 %. Überarbeitet nach /lp-pruefung (07.10.2026):
% neue Teilaufgaben (G1c Ablesen, G2c Glas mit 3 · α, G3c Dichte, G4d Temperatur, G5b Volumen bei −100 °C,
% G5d Manometer), Gleichwertigkeiten auch in Abschnitt 3, Zeilenabstand je Zelle (\lineskip in der X-Spalte).
% Runde 2 (08.10.2026): G2 Messkolben, G3 1750/750 m bei 2 K, G1c mit l = l0 + Δl, Celsius-Fehler in G4 nur einmal
% abgezogen (auch über d), c) folgerichtig aus dem eigenen b), Kelvin als T [K] = ϑ [°C] + 273.15, \addlinespace 6 pt.
\documentclass[11pt]{article}
\usepackage{../lp-druck}
\renewcommand{\lpname}{Wärmeausdehnung und Gase}
\renewcommand{\lpteilgebiet}{5.3}
\renewcommand{\lpfuss}{Bewertungspaket}
\newcolumntype{P}{>{\centering\arraybackslash}p{10mm}}
\newenvironment{raster}{\par\smallskip\noindent\renewcommand{\arraystretch}{1.5}\setlength{\lineskiplimit}{4pt}\setlength{\lineskip}{8pt}\tabularx{\linewidth}{@{}>{\raggedright\arraybackslash}p{38mm}P >{\raggedright\arraybackslash\setlength{\lineskiplimit}{3pt}\setlength{\lineskip}{7pt}}X@{}}\toprule
  \small\textsc{Teilschritt} & \small\textsc{P} & \small\textsc{Musterlösung}\\\midrule}{\bottomrule\endtabularx\par}
\newcommand{\fehler}[1]{\par\smallskip{\small\raggedright\textcolor{lprot}{\bfseries Typische Fehler:} #1\par}}
\newcommand{\E}{\,{\footnotesize\bfseries\color{lpgruen}(E)}}
\newcommand{\B}{\,{\footnotesize\bfseries(B)}}
\newcommand{\W}{\,{\footnotesize\bfseries(W)}}
\newcommand{\A}{\,{\footnotesize\bfseries\color{lpgruen}(A)}}
\newcommand{\pK}{\;\tfrac{1}{\text{K}}}
\begin{document}
\lpkopf{Bewertungspaket zum Gesamttest}{}

\begin{lpachtung}{Erst lösen, dann öffnen}
Dieses Paket enthält die Musterlösung. Wer es vor dem Lösen liest, misst am Ende nur, wie gut das Abschreiben gelingt.
\end{lpachtung}

\section*{1 · So gehst du vor}
\begin{enumerate}[leftmargin=*,itemsep=2pt]
\item \textbf{Gesamttest lösen} — vollständig, mit Rechenweg, ohne dieses Paket.
\item \textbf{Fotografieren oder scannen}, gut lesbar, eine Seite pro Bild. \textbf{Kein Name} auf den Bildern.
  Handyfotos tragen oft unsichtbar Standort, Zeit und Gerät mit (Metadaten). Lade darum lieber einen Scan oder ein
  Bildschirmfoto des Fotos hoch, oder entferne vorher die Standortangaben.
\item \textbf{Bewerten lassen.} Ob du dafür eine KI nutzt und welche, entscheidest du selbst und in eigener Verantwortung —
  je nachdem, was dir aufgrund deines Alters und deiner persönlichen Situation erlaubt ist (Altersgrenzen und
  Nutzungsbedingungen des Anbieters, allenfalls Einverständnis deiner Eltern). Lade nur deine Lösung und dieses PDF hoch,
  keine persönlichen Daten, und füge den Auftrag aus Abschnitt~2 ein. Ohne KI kannst du dich mit Abschnitt~3 selbst bewerten.
\item \textbf{Rückmeldung prüfen.} Stimmt, was die KI aus deiner Handschrift gelesen hat? Eine KI kann sich irren —
  im Zweifel gilt das Raster in Abschnitt~3.
\item \textbf{Gezielt wiederholen} nach Abschnitt~4.
\end{enumerate}

\needspace{8\baselineskip}
\section*{2 · Auftrag an die KI}
\begin{tcolorbox}[breakable,colback=lplinie!25,colframe=lplinie,boxrule=0.5pt,arc=1mm,fontupper=\small]
Du bist Korrektorin oder Korrektor für Physik an einer Schweizer Berufsmaturitätsschule. Im Anhang findest du (1)~das Bewertungspaket zum Gesamttest «Wärmeausdehnung und Gase» mit Musterlösung und Punkteraster und (2)~Fotos meiner handschriftlichen Lösung.\par\medskip
Geh so vor:\par
1. Schreib zu jeder Aufgabe G1 bis G5 zuerst ab, was du in meiner Lösung liest — Rechenweg und Ergebnis. Was du nicht sicher lesen kannst, markierst du mit [unsicher]. Nicht raten.\par
2. Vergib die Punkte genau nach dem Raster im Bewertungspaket (Abschnitt 3), ganze Punkte je Zeile. Ziehe einen Fehler nur einmal ab: Rechne ich mit einem falschen Zwischenergebnis richtig weiter, gibt es die Folgepunkte, und derselbe Fehler (etwa dieselbe fehlende Umrechnung) in mehreren Teilaufgaben kostet nur einen Punkt — auch wenn das Folgeergebnis unplausibel ist; eine fehlende Plausibilitätsprobe kostet keinen zusätzlichen Punkt. Die Zeilen «Typische Fehler» sagen, welche Rasterzeile bei diesem Fehler 0 Punkte gibt — sie sind kein zusätzlicher Abzug.\par
3. Andere richtige Lösungswege zählen voll. Gleichwertige Schreibweisen zählen voll: Dezimalpunkt oder Dezimalkomma, andere Einheit oder Vorsilbe ($45\;\text{mm} = 0.045\;\text{m}$; $1.1\;\text{cm}^3 = 1100\;\text{mm}^3$; $1.26\;\text{l} = 1260\;\text{cm}^3$; $l_0 = 125\,000\;\text{mm}$ statt $125\;\text{m}$; $h_0$ aus cm-Werten in cm gerechnet und richtig in m umgewandelt; $1/\text{K} = \text{K}^{-1}$), Zehnerpotenz oder ausgeschriebene Zahl, sinnvoll gerundete Werte (drei signifikante Stellen), in $T\;[\text{K}] = \vartheta\;[^\circ\text{C}] + 273.15$ auch $+ 273$ statt $+ 273.15$. Zeilen mit (E) sind Ergebnispunkte: Sie verlangen das Ergebnis \emph{mit richtiger Einheit}, auch ohne Weg; ohne oder mit falscher Einheit gibt es diesen Punkt nicht. Zeilen mit (A) sind Ablesepunkte: Es zählt der abgelesene Wert mit Einheit, innerhalb der angegebenen Ablesetoleranz. Zeilen mit (W) gibt es nur mit nachvollziehbarem Weg (Formel, dann Werte mit Einheiten); ein Zahlenwert darin braucht ebenfalls die Einheit. Zeilen mit (B) sind Begründungen: Es zählt die fachliche Aussage in eigenen Worten, eine Formel ist nicht nötig; andere richtige Formulierungen zählen voll. Werte aus sinnvoll gerundeten Zwischenergebnissen zählen, wenn sie höchstens rund 2\,\% vom Raster abweichen.\par
4. Begründe jeden Abzug in einem Satz und nenne die Fehlerart (zum Beispiel «Temperatur in Grad Celsius statt in Kelvin»). Schreib die Musterlösung nicht ab — zeig mir, wo mein Fehler liegt.\par
5. Gib zum Schluss eine Tabelle «Aufgabe | Punkte | Begründung», die Gesamtpunktzahl (von 25) und — nach der Selbsteinschätzung in Abschnitt 4 — welches Kapitel des Leitprogramms ich wiederholen soll. Keine Note.\par
6. Fehlt ein Foto oder ist eine Aufgabe unlesbar, sag es, statt sie zu bewerten.
\end{tcolorbox}

\section*{3 · Musterlösung und Punkteraster}
\textbf{Allgemein:} Ganze Punkte je Zeile. Ein Fehler wird nur einmal abgezogen: Wer mit einem falschen Zwischenergebnis
richtig weiterrechnet, bekommt die folgenden Zeilen (Folgefehler), und derselbe Fehler in mehreren Teilaufgaben kostet nur einen Punkt — auch wenn das Folgeergebnis unplausibel ist; eine fehlende
Plausibilitätsprobe kostet keinen zusätzlichen Punkt. Gleichwertige Wege und Schreibweisen zählen voll: Dezimalpunkt oder
Dezimalkomma, andere Einheit oder Vorsilbe (\(45\;\text{mm} = 0.045\;\text{m}\); \(l_0 = 125\,000\;\text{mm}\) ist richtig,
falsch ist nur \(125\,000\;\text{m}\)), Zehnerpotenz oder ausgeschriebene Zahl, sinnvoll gerundete Werte, in \(T\;[\text{K}] = \vartheta\;[^\circ\text{C}] + 273.15\) auch \(+ 273\) statt \(+ 273.15\).
In der Spalte P: \textbf{(E)}~= Ergebnispunkt — es zählt das Ergebnis \emph{mit richtiger Einheit}, auch ohne Weg; ohne oder mit
falscher Einheit gibt es ihn nicht. \textbf{(A)}~= Ablesepunkt — der abgelesene Wert mit Einheit, innerhalb der angegebenen Toleranz. \textbf{(W)}~= Wegpunkt — nur mit nachvollziehbarem Weg (Formel, dann Werte mit Einheiten);
steht in einer Weg-Zeile ein Zahlenwert, braucht auch er die Einheit. \textbf{(B)}~= Begründungspunkt — es zählt die fachliche
Aussage in eigenen Worten, eine Formel ist nicht nötig. Folgewerte aus gerundeten Zwischenergebnissen zählen, wenn sie höchstens rund 2\,\% abweichen.
«Typische Fehler» nennt die Zeilen mit 0 Punkten und die Punktzahl, die bleibt.

\begin{lpaufgabe}{G1}{Eine Brücke über das Jahr}{5 P}
\begin{raster}
a) Ablesen & 1\A & Einbau bei \(10\;^\circ\text{C}\) (dort ist \(\Delta l = 0\)); bei \(40\;^\circ\text{C}\) ist sie \(45\;\text{mm}\) länger (zählt: \(\pm 1\;^\circ\text{C}\) und \(\pm 2\;\text{mm}\); beide Werte nötig)\\\addlinespace[6pt]
b) Ansatz & 1\W & \(\Delta l = \alpha \cdot l_0 \cdot \Delta T\), also \(l_0 = \dfrac{\Delta l}{\alpha \cdot \Delta T}\) mit \(\Delta T = 40\;^\circ\text{C} - 10\;^\circ\text{C} = 30\;\text{K}\)\\\addlinespace[6pt]
b) Länge & 1\E & \(l_0 = \dfrac{0.045\;\text{m}}{12 \cdot 10^{-6}\pK \cdot 30\;\text{K}} = 125\;\text{m}\) (folgerichtig aus den eigenen Ablesewerten; \(125\,000\;\text{mm}\) zählt)\\\addlinespace[6pt]
c) Temperatur, Länge & 1\E & \(\Delta l = -30\;\text{mm}\) bei \(-10\;^\circ\text{C}\) (zählt: \(\pm 1\;^\circ\text{C}\)); \(l = l_0 + \Delta l = 125\;\text{m} - 0.030\;\text{m} = 124.97\;\text{m}\) bzw. \(124\,970\;\text{mm}\) (folgerichtig aus b). \(125\;\text{m}\) zählt hier nicht (Ausnahme von der Rundungsregel). Beides nötig\\\addlinespace[6pt]
d) Aluminium & 1\B & Steiler, rund doppelt so steil, weil \(\alpha\) knapp doppelt so gross ist; ebenfalls durch \(\Delta l = 0\) bei \(10\;^\circ\text{C}\) (bei \(40\;^\circ\text{C}\) rund \(89\;\text{mm}\)). Nötig ist «steiler» mit Begründung über \(\alpha\); der Punkt bei \(10\;^\circ\text{C}\) ist erwünscht, aber nicht nötig. Nur «steiler» ohne \(\alpha\): 0\\\addlinespace[6pt]
\end{raster}
\fehler{Temperatur statt Differenz (\(\Delta T = 40\;\text{K}\), \(l_0 \approx 93.8\;\text{m}\)): b) Ansatz 0, Länge folgerichtig \(\to\) 4 von 5\,P. mm nicht in m umgerechnet (\(l_0 = 125\,000\;\text{m}\)): b) Länge 0 \(\to\) 4 von 5\,P. In c) \(+30\;^\circ\text{C}\) (länger statt kürzer gelesen): c) 0 \(\to\) 4 von 5\,P.}
\end{lpaufgabe}

\begin{lpaufgabe}{G2}{Der Messkolben}{5 P}
\begin{raster}
a) Ansatz & 1\W & \(\Delta V = \gamma \cdot V_0 \cdot \Delta T\) mit \(\Delta T = 24\;^\circ\text{C} - 20\;^\circ\text{C} = 4\;\text{K}\) (Flüssigkeit: \(\gamma\) aus der Angabe, kein Faktor 3)\\\addlinespace[6pt]
a) Volumen & 1\E & \(\Delta V = 1.10 \cdot 10^{-3}\pK \cdot 250\;\text{cm}^3 \cdot 4\;\text{K} = 1.1\;\text{cm}^3\) (\(= 1100\;\text{mm}^3\))\\\addlinespace[6pt]
b) Pegel & 1\E & Säule mit \(V = A \cdot h\): \(h = \dfrac{\Delta V}{A} = \dfrac{1.1\;\text{cm}^3}{0.80\;\text{cm}^2} \approx 1.4\;\text{cm} = 14\;\text{mm}\) (folgerichtig aus a; \(13.75\;\text{mm}\) zählt)\\\addlinespace[6pt]
c) Glas & 1\E & Festkörper: \(\gamma_\text{Glas} \approx 3 \cdot \alpha_\text{Glas} = 24.3 \cdot 10^{-6}\pK\); \(\Delta V_\text{Glas} = 24.3 \cdot 10^{-6}\pK \cdot 250\;\text{cm}^3 \cdot 4\;\text{K} \approx 0.024\;\text{cm}^3\) (\(= 24\;\text{mm}^3\))\\\addlinespace[6pt]
d) mehr/weniger & 1\B & Etwas weniger: Der Innenraum wird auch grösser; rund \(0.024\;\text{cm}^3\) des Zuwachses bleiben im Kolben, nur der Rest (rund \(1.08\;\text{cm}^3\)) steigt in den Hals (rund \(13.4\;\text{mm}\)). Der Unterschied ist klein (rund \(2\,\%\)). «Weniger» mit Begründung nötig\\\addlinespace[6pt]
\end{raster}
\fehler{Mit \(3 \cdot \gamma\) beim Ethanol gerechnet (\(3.3\;\text{cm}^3\)): a) Volumen 0, b) folgerichtig (\(41\;\text{mm}\)) \(\to\) 4 von 5\,P. cm nicht in mm umgerechnet und \(1.375\;\text{mm}\) angegeben: b) 0 \(\to\) 4 von 5\,P. In c) \(\alpha\) statt \(3 \cdot \alpha\) (\(0.0081\;\text{cm}^3\)): c) 0 \(\to\) 4 von 5\,P; d) zählt weiter. In d) «mehr, weil das Glas mitdrückt»: d) 0.}
\end{lpaufgabe}

\begin{lpaufgabe}{G3}{Zwei Annahmen für ein Meer}{5 P}
\begin{raster}
a) Ablesen & 1\A & A: \(73.5\;\text{cm}\), B: \(31.5\;\text{cm}\) bei \(\Delta T = 2\;\text{K}\) (zählt: je \(\pm 2\;\text{cm}\); beide Werte nötig)\\\addlinespace[6pt]
b) Ansatz & 1\W & \(\Delta h = \gamma \cdot h_0 \cdot \Delta T\), also \(h_0 = \dfrac{\Delta h}{\gamma \cdot \Delta T}\)\\\addlinespace[6pt]
b) Schichtdicken & 1\E & A: \(h_0 = \dfrac{0.735\;\text{m}}{0.21 \cdot 10^{-3}\pK \cdot 2\;\text{K}} = 1750\;\text{m}\); B: \(h_0 = \dfrac{0.315\;\text{m}}{0.21 \cdot 10^{-3}\pK \cdot 2\;\text{K}} = 750\;\text{m}\) (beide nötig; folgerichtig aus den Ablesewerten; in cm gerechnet und \(175\,000\;\text{cm}\) usw. angegeben zählt)\\\addlinespace[6pt]
c) Dichte & 1\B & \(\rho = \dfrac{m}{V}\): Die Masse bleibt gleich, das Volumen wächst beim Erwärmen, also sinkt die Dichte. Das wärmere, weniger dichte Wasser schwimmt auf dem dichteren, kälteren (Auftrieb). Beide Teile nötig\\\addlinespace[6pt]
d) Wasserthermometer & 1\B & Wasser zieht sich von \(0\;^\circ\text{C}\) bis \(4\;^\circ\text{C}\) beim Erwärmen zusammen und dehnt sich danach wieder aus. Derselbe Fadenstand gehört dann zu zwei Temperaturen (etwa \(0\;^\circ\text{C}\) und \(8\;^\circ\text{C}\)) — die Anzeige ist nicht eindeutig\\\addlinespace[6pt]
\end{raster}
\fehler{cm nicht in m umgerechnet (\(h_0 = 175\,000\;\text{m}\) und \(75\,000\;\text{m}\)): derselbe Fehler bei A und B, einmal abgezogen: b) Schichtdicken 0 \(\to\) 4 von 5\,P. Mit \(\Delta T = 1\;\text{K}\) statt \(2\;\text{K}\) (\(3500\;\text{m}\) und \(1500\;\text{m}\)): b) Schichtdicken 0 \(\to\) 4 von 5\,P. In c) «warmes Wasser ist leichter» ohne Volumen oder Dichte: c) 0. Hinweis zu c): Wer ergänzt, dass Meerwasser (mit Salz) keine Anomalie bei \(4\;^\circ\text{C}\) hat, liegt richtig; die Frage verlangt es nicht.}
\end{lpaufgabe}

\begin{lpaufgabe}{G4}{Die PET-Flasche vom Pass}{5 P}
\begin{raster}
a) Ansatz & 1\W & \(\dfrac{p_1 \cdot V_1}{T_1} = \dfrac{p_2 \cdot V_2}{T_2}\), also \(V_2 = \dfrac{p_1 \cdot V_1 \cdot T_2}{T_1 \cdot p_2}\), Temperaturen in Kelvin\\\addlinespace[6pt]
a) Volumen & 1\E & \(V_2 = \dfrac{0.75\;\text{bar} \cdot 1.5\;\text{l} \cdot 298.15\;\text{K}}{278.15\;\text{K} \cdot 0.96\;\text{bar}} \approx 1.26\;\text{l}\)\\\addlinespace[6pt]
b) starr & 1\E & Volumen gleich: \(p_2 = p_1 \cdot \dfrac{T_2}{T_1} = 0.75\;\text{bar} \cdot \dfrac{298.15\;\text{K}}{278.15\;\text{K}} \approx 0.80\;\text{bar}\)\\\addlinespace[6pt]
c) Begründung & 1\B & Innen wären es nur rund \(0.80\;\text{bar}\), aussen drücken \(0.96\;\text{bar}\): Der grössere Aussendruck drückt die weiche Flasche ein, bis innen und aussen gleich viel Druck herrscht (bei rund \(1.26\;\text{l}\)). Nach einem Folgefehler in b) zählt c) folgerichtig: Es zählt der richtige Vergleich des eigenen Werts aus b) mit dem Aussendruck, auch wenn er dann nicht «eingedrückt» ergibt\\\addlinespace[6pt]
d) Temperatur & 1\E & \(T_2 = \dfrac{p_2 \cdot V_2 \cdot T_1}{p_1 \cdot V_1} = \dfrac{0.96\;\text{bar} \cdot 1.5\;\text{l} \cdot 278.15\;\text{K}}{0.75\;\text{bar} \cdot 1.5\;\text{l}} \approx 356\;\text{K}\), also rund \(83\;^\circ\text{C}\) (nur in Kelvin angegeben: 0)\\\addlinespace[6pt]
\end{raster}
\fehler{Drücke vertauscht (\(V_2 \approx 2.06\;\text{l}\)): a) Volumen 0 \(\to\) 4 von 5\,P. Durchgehend in Grad Celsius gerechnet (\(V_2 \approx 5.9\;\text{l}\), \(p \approx 3.75\;\text{bar}\), in d) \(6.4\;^\circ\text{C}\)): derselbe Fehler, nur einmal abgezogen: a) Ansatz 0; a) Volumen, b), c) (Regel oben) und d) folgerichtig \(\to\) 4 von 5\,P. Nur in d) in Grad Celsius gerechnet, a) und b) richtig: d) 0 \(\to\) 4 von 5\,P. In c) nur «weil es wärmer wird»: c) 0.}
\end{lpaufgabe}

\begin{lpaufgabe}{G5}{Eine Gerade in Grad Celsius}{5 P}
\begin{raster}
a) Ablesen & 1\A & \(2.0\;\text{l}\) bei \(0\;^\circ\text{C}\), \(3.1\;\text{l}\) bei \(150\;^\circ\text{C}\) (zählt: je \(\pm 0.05\;\text{l}\); beide nötig)\\\addlinespace[6pt]
b) Ansatz & 1\W & Isobar: \(\dfrac{V_1}{T_1} = \dfrac{V_2}{T_2}\) mit \(T_1 = 273.15\;\text{K}\), \(T_2\;[\text{K}] = -100 + 273.15 = 173.15\)\\\addlinespace[6pt]
b) Volumen & 1\E & \(V_2 = V_1 \cdot \dfrac{T_2}{T_1} = 2.0\;\text{l} \cdot \dfrac{173.15\;\text{K}}{273.15\;\text{K}} \approx 1.27\;\text{l}\) (folgerichtig aus dem eigenen Ablesewert; vom Punkt bei \(150\;^\circ\text{C}\) aus gleichwertig)\\\addlinespace[6pt]
c) Grenze & 1\B & Ein Grund genügt: Das Gas wird vorher flüssig oder fest, dann gilt das Modell des idealen Gases nicht mehr; die Teilchen haben ein Eigenvolumen; \(0\;\text{K}\) (\(-273.15\;^\circ\text{C}\)) ist nicht erreichbar\\\addlinespace[6pt]
d) Manometer & 1\E & Vorher absolut \(p_1 = 1.0\;\text{bar}\) (Manometer \(0\;\text{bar}\) + Luftdruck). Isotherm: \(p_2 = \dfrac{p_1 \cdot V_1}{V_2} = \dfrac{1.0\;\text{bar} \cdot 2.0\;\text{l}}{1.6\;\text{l}} = 1.25\;\text{bar}\) absolut; das Manometer zeigt \(1.25\;\text{bar} - 1.0\;\text{bar} = 0.25\;\text{bar}\) (folgerichtig mit dem eigenen Ablesewert bei \(0\;^\circ\text{C}\))\\\addlinespace[6pt]
\end{raster}
\fehler{In b) mit Grad Celsius gerechnet (vom Punkt bei \(150\;^\circ\text{C}\) aus \(V \approx -2.07\;\text{l}\)): b) Ansatz 0, Volumen folgerichtig \(\to\) 4 von 5\,P. In d) \(1.25\;\text{bar}\) als Anzeige des Manometers: d) 0 \(\to\) 4 von 5\,P. In d) mit \(p_1 = 0\;\text{bar}\) gerechnet: d) 0. Druck proportional statt umgekehrt (\(0.8\;\text{bar}\) absolut): d) 0.}
\end{lpaufgabe}

\needspace{10\baselineskip}
\section*{4 · Selbsteinschätzung}
\begin{tabularx}{\linewidth}{@{}l X@{}}\toprule
22 – 25 P & Die geprüften Teile sitzen. Wo du Punkte verloren hast: das Kapitel dieser Aufgabe nochmals (Zuordnung unten).\\
17 – 21 P & Den schwächsten Teil nochmals: Simulation und Übungen des Kapitels, in dem du die meisten Punkte verloren hast.\\
11 – 16 P & Zurück zu den Kapiteln aller Aufgaben, in denen du Punkte verloren hast.\\
0 – 10 P & Zurück zu Kapitel 1 und von dort der Reihe nach weiter.\\\midrule
\multicolumn{2}{@{}p{\linewidth}@{}}{\small\textbf{Aufgabe $\to$ Kapitel} (gilt für alle Bereiche): G1 $\to$ Kapitel 1 (K1) · G2 $\to$ Kapitel 2 (K1) · G3 $\to$ Kapitel 3 (K1) · G4 $\to$ Kapitel 4 (K2) · G5 $\to$ Kapitel 5 (K2), G5d (Manometer, absoluter Druck) auch Kapitel 4}\\\bottomrule
\end{tabularx}
\end{document}
